\begin{trapbox}

	\begin{description}[style=nextline, leftmargin=2.8em, labelsep=0.5em, itemsep=0.35em]
		\item[\textbf{\textcolor{CTrap}{易错点一（分母非零条件）：}}]
			构造 $b_n=\dfrac{a_n-\alpha}{a_n-\beta}$ 时，要求分母 $C a_n+D\neq0$；若某 $a_n$ 使 $C a_n+D=0$，则递推无定义，公式失效。
		\item[\textbf{\textcolor{CTrap}{正确理解（分母需检验）：}}]
			使用不动点法前应先确认递推式对所有项都有意义，初项应使分母非零，并在递推过程中保持分母非零。
		\item[\textbf{\textcolor{CTrap}{错因分析（盲目套公式）：}}]
			学生直接套用定理，忽略推导中分母非零的假设，导致在边界项出错。
	\end{description}

	\medskip
	\begin{description}[style=nextline, leftmargin=2.8em, labelsep=0.5em, itemsep=0.35em]
		\item[\textbf{\textcolor{CTrap}{易错点二（重根误用等比）：}}]
			当不动点 $\alpha=\beta$ 时，若仍使用比例公式，则会得到 $b_{n+1}=b_n$ 的恒等式，无法求出通项。
		\item[\textbf{\textcolor{CTrap}{正确理解（重根取倒数）：}}]
			重根时应取倒数 $c_n=\dfrac{1}{a_n-\alpha}$，得到等差数列而非等比数列。
		\item[\textbf{\textcolor{CTrap}{错因分析（结论混淆）：}}]
			学生未判断根的情况，习惯性使用比例法，忽视了重根的特性。
	\end{description}

	\medskip
	\begin{description}[style=nextline, leftmargin=2.8em, labelsep=0.5em, itemsep=0.35em]
		\item[\textbf{\textcolor{CTrap}{易错点三（系数退化情形）：}}]
			当 $A-\alpha C=0$ 或 $A-\beta C=0$ 时，比例公式中分母或分子为零，公式不成立；此时数列可能退化为常数列等简单情形。
		\item[\textbf{\textcolor{CTrap}{正确理解（条件再确认）：}}]
			使用等比公式前必须验证 $A-\alpha C\neq0$ 且 $A-\beta C\neq0$；若不满足，应直接分析递推或利用定义求通项。
		\item[\textbf{\textcolor{CTrap}{错因分析（忽视前提）：}}]
			学生只记忆比例形式，未关注定理中给出的非零条件，导致公式滥用。
	\end{description}

\end{trapbox}
